% !TeX root = ./../main.tex

\section{Uczenie maszynowe}\label{sec:tb:ml}
\agentinfo{
Stan (PHD-TS-95, 2026-11-27): podrozdział skrócony i poprawiony po recenzji \agentref{eksperymenty/recenzje/rozdzial-2-sekcje-2.5-2.6.md}; usunięto rysunek modelu pix2pix, miary klasyfikacji poza dokładnością, własności gry GAN, jądro kwantowe, wzór reguły przesunięcia parametru i przegląd sieci ekwiwariantnych; dopisano akapit o porównaniach sparowanych i wnioskowaniu statystycznym (etykieta par:tb:ml:stats) oraz gradient z modelu zastępczego; lista zmian w \agentref{eksperymenty/recenzje/rozdzial-2-skracanie.md}.
Zakres: wyłącznie pojęcia używane w rozdziałach 4 i 6--8 oraz w sekcjach 3.5 i 3.7; podrozdział 2.6.5 łączy dawne podrozdziały 2.1.6 i 2.6.5 (etykiety subsec:tb:lie, subsec:tb:ml:lie, eq:tb01:generator, eq:tb:ml:lie) i służy tylko akapitom par:lr:noise:generators i par:lr:cls:lie oraz zdaniu o kwantowych sieciach ekwiwariantnych w sekcji 3.7.
Architektury, hiperparametry i liczby z eksperymentów opisują rozdziały 6--8; tutaj tylko definicje.
Miary jakości obrazu i współczynniki korelacji definiuje podrozdział o reprezentacjach obrazów.
Notacja dla rozdziałów 6--8: generator i dyskryminator to makra netG i netD, sieć wyznaczająca kąty obwodu to Theta z indeksem w, liczba pomiarów to makro nshots.
Etykiety używane przez inne pliki (nie usuwać): etykiety podsekcji (w tym subsec:tb:lie), eq:tb:ml:metrics, eq:tb:ml:adam, fig:tb:ml:hybrid, eq:tb:ml:lie, eq:tb01:generator.
Akapit o wnioskowaniu statystycznym nie ma cytowań dla testu t, poprawki Holma i przekształcenia Fishera (brak zweryfikowanych wpisów w bazie); rozdziały 4, 7 i 8 nie odwołują się jeszcze do etykiety par:tb:ml:stats.
}

Podrozdział wprowadza tylko pojęcia używane w dalszych rozdziałach; podstawy wykłada monografia~\cite{goodfellow2016deep}.

\subsection{Uczenie nadzorowane i klasyfikacja obrazów}\label{subsec:tb:ml:supervised}

W uczeniu nadzorowanym dany jest zbiór danych $\dataset=\{(\mathbf{x}_i,\mathbf{y}_i)\}_{i}$, w którym każda próbka $\mathbf{x}_i$ (np. obraz jako wektor intensywności pikseli) ma przypisaną wartość docelową $\mathbf{y}_i$.
W klasyfikacji jest nią etykieta jednej z $K$ klas, w translacji obraz--obraz z rozdziału~\ref{chapter:dl_error_reduction} obraz odniesienia, a w autoenkoderze z rozdziału~\ref{chapter:learned_encoding} sam obraz wejściowy.
Model $f_{\mathbf{w}}$ o parametrach (wagach) $\mathbf{w}$ wyznacza predykcję $\hat{\mathbf{y}}=f_{\mathbf{w}}(\mathbf{x})$, a uczenie polega na minimalizacji średniej straty na zbiorze uczącym:
\begin{equation}\label{eq:tb:ml:risk}
\symloss(\mathbf{w})=\frac{1}{|\trainset|}\sum_{(\mathbf{x},\mathbf{y})\in\trainset}\ell\bigl(f_{\mathbf{w}}(\mathbf{x}),\mathbf{y}\bigr).
\end{equation}
Dla wartości ciągłych stratą $\ell$ jest średni błąd bezwzględny (MAE), nazywany wtedy stratą L1 i oznaczany $\ell_{\mathrm{L1}}$, albo błąd średniokwadratowy~\eqref{eq:tb:qimg:errors}.
Dla klasyfikatora zwracającego rozkład prawdopodobieństwa klas $\hat{\mathbf{y}}$ stratą jest entropia krzyżowa $-\sum_{k}y_k\ln\hat{y}_k$, w której $\mathbf{y}$ jest etykietą zapisaną jako wektor jedynkowy, a dla jednego wyjścia $\hat{y}\in(0,1)$ i etykiety $y\in\{0,1\}$ binarna entropia krzyżowa (BCE) $-y\ln\hat{y}-(1-y)\ln(1-\hat{y})$.

Dane dzieli się na rozłączne zbiory: uczący $\trainset$, walidacyjny $\valset$, na którym dobiera się hiperparametry i wybiera punkt kontrolny modelu (zapisany stan wag), oraz testowy $\testset$, używany jednokrotnie do oceny końcowej, ponieważ wybór najlepszego z wielu wariantów na podstawie wyniku na danym zbiorze zawyża ten wynik.
\Gls{overfitting} zachodzi, gdy model dopasowuje się do szczegółów i szumu zbioru uczącego: strata na $\trainset$ maleje, a na $\valset$ rośnie; ryzyko to jest tym większe, im większa jest pojemność modelu względem liczby próbek~\cite{goodfellow2016deep}.
Od przeuczenia trzeba odróżnić \glsdisp{data-leakage}{przeciek danych} (ang.~\emph{data leakage}): wpływ informacji ze zbioru ocenianego na uczenie lub wybór modelu, który zawyża ocenę niezależnie od rzeczywistej zdolności uogólniania~\cite{kapoor_leakage_2023}.
W rozdziałach~\ref{chapter:learned_encoding} i~\ref{chapter:classification} znaczenie mają trzy jego formy: te same lub niemal te same próbki w zbiorze uczącym i ocenianym, parametry przetwarzania wstępnego wyznaczone przed podziałem danych oraz dobór hiperparametrów lub punktu kontrolnego na zbiorze, na którym podaje się wynik.

Wynik klasyfikatora podsumowuje macierz pomyłek $\mathbf{C}$, której element $C_{kl}$ jest liczbą próbek klasy $k$ przypisanych do klasy $l$, i wyznaczana z niej dokładność
\begin{equation}\label{eq:tb:ml:metrics}
\mathrm{ACC}=\frac{\sum_{k}C_{kk}}{\sum_{k,l}C_{kl}}.
\end{equation}
Dokładność odnosi się do klasyfikatora trywialnego: przy $K$ równolicznych klasach zgadywanie daje $1/K$.

\paragraph{Porównania sparowane i wnioskowanie statystyczne.}\label{par:tb:ml:stats}
Dwie metody ocenione na tych samych jednostkach (obrazach, próbkach, ziarnach generatora liczb losowych) porównuje się parami, przez różnice wyników.
Przedział ufności średniej różnicy wyznacza się z rozkładu t~Studenta o liczbie stopni swobody o jeden mniejszej od liczby par (sparowany test t sprawdza wtedy hipotezę o zerowej średniej różnicy i zakłada jej rozkład normalny) albo metodą bootstrap~\cite{efron_bootstrap_1979}, czyli z rozkładu średnich obliczonych na próbach losowanych ze zwracaniem spośród zaobserwowanych różnic.
Testy nieparametryczne nie zakładają rozkładu normalnego~\cite{bib:statistics_wilc}: dokładny test znaków korzysta tylko ze znaków różnic, które przy braku różnicy między metodami mają rozkład dwumianowy z prawdopodobieństwem $1/2$, test Wilcoxona także z rang ich modułów, test permutacyjny wyznacza rozkład statystyki przez losowe zamiany metod w parach, a test McNemara~\cite{mcnemar_note_1947} porównuje dwa klasyfikatory na podstawie liczby próbek, które poprawnie rozpoznał tylko jeden z nich.
W rodzinie $k$ testów stosuje się poprawkę Holma: p-wartości uporządkowane rosnąco porównuje się kolejno z progami $\alpha/k$, $\alpha/(k-1)$, \dots, do pierwszej hipotezy, której nie można odrzucić.
Przedział ufności współczynnika korelacji $r$ z $N$ obserwacji wyznacza się przekształceniem Fishera ($\operatorname{artanh}r$ ma w przybliżeniu rozkład normalny o odchyleniu standardowym $1/\sqrt{N-3}$), a współczynnik determinacji $R^2$ jest udziałem wariancji danych wyjaśnionym przez dopasowany model.

\subsection{Sieci neuronowe i sieci splotowe}\label{subsec:tb:ml:cnn}

Neuron oblicza $h_j=a\bigl(\sum_{i}w_{ji}x_i+b_j\bigr)$, gdzie $w_{ji}$ są wagami, $b_j$ przesunięciem, a $a$ nieliniową funkcją aktywacji; warstwy neuronów złożone kolejno tworzą perceptron wielowarstwowy (MLP).
W pracy występują funkcje aktywacji $\mathrm{ReLU}(u)=\max(0,u)$, $\mathrm{LeakyReLU}$, sigmoidalna $1/(1+e^{-u})$, $\tanh$ oraz softmax, zamieniająca wektor na rozkład prawdopodobieństwa.
Gradient straty~\eqref{eq:tb:ml:risk} oblicza się propagacją wsteczną~\cite{rumelhart_learning_1986}, szczególnym przypadkiem automatycznego różniczkowania~\cite{goodfellow2016deep}, i szacuje na losowych paczkach (wsadach) próbek; przejście przez cały zbiór uczący nazywa się epoką.
Optymalizator Adam~\cite{kingma2014adam} aktualizuje wagi na podstawie wygładzonych wykładniczo estymat pierwszego i drugiego momentu gradientu $\mathbf{g}_t$:
\begin{equation}\label{eq:tb:ml:adam}
\mathbf{m}_t=\beta_1\mathbf{m}_{t-1}+(1-\beta_1)\mathbf{g}_t,\quad
\mathbf{v}_t=\beta_2\mathbf{v}_{t-1}+(1-\beta_2)\mathbf{g}_t^{2},\quad
\mathbf{w}_t=\mathbf{w}_{t-1}-\eta\,\frac{\hat{\mathbf{m}}_t}{\sqrt{\hat{\mathbf{v}}_t}+\epsilon},
\end{equation}
gdzie działania wykonuje się po składowych, daszek oznacza korektę obciążenia ($\hat{\mathbf{m}}_t=\mathbf{m}_t/(1-\beta_1^{t})$, podobnie $\hat{\mathbf{v}}_t$), $\eta$ jest współczynnikiem uczenia, $\beta_1,\beta_2\in[0,1)$ współczynnikami zaniku, a $\epsilon$ małą stałą.

W warstwie splotowej sieci splotowej (CNN)~\cite{lecun1998gradient} każdy neuron oblicza ważoną sumę wartości z okna wokół danego położenia, przy czym te same wagi (jądro splotu) są używane we wszystkich położeniach~\cite{goodfellow2016deep}; splot z krokiem 2 zmniejsza rozdzielczość mapy cech dwukrotnie, a splot transponowany ją zwiększa.
Normalizacja wsadowa~\cite{ioffe_batch_2015} standaryzuje aktywacje średnią i wariancją z bieżącej paczki, a dropout~\cite{srivastava_dropout_2014} losowo usuwa podczas uczenia neurony wraz z ich połączeniami, co ogranicza przeuczenie.
Sieć U-Net~\cite{ronneberger_u-net_2015} składa się ze ścieżki kodującej, w której rozdzielczość map cech maleje, i ścieżki dekodującej, do której poziomów dołącza się mapy cech z odpowiednich poziomów kodera (połączenia skrótowe)~\cite{isola_image_2017}.

\subsection{Generatywne sieci przeciwstawne}\label{subsec:tb:ml:gan}

Generatywna sieć przeciwstawna (GAN)~\cite{goodfellow2014generative} składa się z generatora $\netG$, który przekształca wektor losowy $\mathbf{z}$ w próbkę, i dyskryminatora $\netD$, który zwraca prawdopodobieństwo, że próbka pochodzi ze zbioru danych, a nie z generatora; w warunkowej sieci GAN obie sieci otrzymują ponadto informację warunkującą $\mathbf{x}$~\cite{mirza2014conditional}.
Sieci uczy się jednocześnie w grze minimaksowej z funkcją wartości
\begin{equation}\label{eq:tb:ml:gan}
V_{\mathrm{c}}(\netD,\netG)=\mathbb{E}_{\mathbf{x},\mathbf{y}}\bigl[\ln\netD(\mathbf{x},\mathbf{y})\bigr]+\mathbb{E}_{\mathbf{x},\mathbf{z}}\bigl[\ln\bigl(1-\netD(\mathbf{x},\netG(\mathbf{x},\mathbf{z}))\bigr)\bigr],
\end{equation}
którą dyskryminator maksymalizuje, czyli minimalizuje stratę BCE z etykietą $1$ dla próbek prawdziwych $\mathbf{y}$ i $0$ dla wygenerowanych.
Model pix2pix~\cite{isola_image_2017} stosuje taką sieć do translacji obraz--obraz: warunkiem jest obraz wejściowy $\mathbf{x}$, generator ma zwrócić obraz docelowy $\mathbf{y}$, a funkcja celu łączy składnik przeciwstawny ze stratą L1 o wadze $\lambda_{\mathrm{L1}}$:
\begin{equation}\label{eq:tb:ml:pix2pix}
\netG^{*}=\arg\min_{\netG}\max_{\netD}\Bigl(V_{\mathrm{c}}(\netD,\netG)+\lambda_{\mathrm{L1}}\,\mathbb{E}_{\mathbf{x},\mathbf{y},\mathbf{z}}\bigl[\ell_{\mathrm{L1}}\bigl(\netG(\mathbf{x},\mathbf{z}),\mathbf{y}\bigr)\bigr]\Bigr).
\end{equation}
Generatorem jest sieć U-Net, a dyskryminator PatchGAN jest siecią w pełni splotową, która ocenia osobno każdy fragment (łatę) obrazu; losowość wprowadza jedynie dropout aktywny również podczas wnioskowania, więc wyniki dla tego samego wejścia mogą się nieznacznie różnić między wywołaniami~\cite{isola_image_2017}.
W eksperymentach z rozdziału~\ref{chapter:dl_error_reduction} i preprintu~\cite{werner_deep_learning_manuscript} strata generatora jest sumą straty BCE dla wyjścia dyskryminatora i błędu MAE względem obrazu odniesienia z wagami $1$ i $100$, co odpowiada $\lambda_{\mathrm{L1}}=100$; obie sieci uczy się optymalizatorem Adam~\eqref{eq:tb:ml:adam}.

\subsection{Kwantowe uczenie maszynowe}\label{subsec:tb:ml:qml}

Kwantowe uczenie maszynowe (QML) obejmuje modele, w których część obliczeń wykonuje obwód kwantowy~\cite{schuld2019qml,benedetti2019vqc}; podstawowym z nich jest wariacyjny obwód kwantowy (VQC), czyli sparametryzowany obwód $U(\boldsymbol{\theta})$ o ustalonej strukturze (ansatz), którego parametry, np. kąty bramek obrotu, dobiera klasyczny optymalizator przez minimalizację funkcji kosztu~\cite{cerezo_variational_2021}.
Dane klasyczne koduje się w stanie kwantowym, np. w kątach obrotów~\cite{benedetti2019vqc}, a wyjściem modelu są wartości oczekiwane obserwabli albo prawdopodobieństwa stanów bazowych, szacowane na urządzeniu z $\nshots$ powtórzeń obwodu.

\Gls{hybrid-model} łączy obwód kwantowy z klasycznym przetwarzaniem wstępnym, które zamienia dane na parametry obwodu, i klasycznym przetwarzaniem końcowym wyników pomiaru; wszystkie te składniki mogą mieć uczone parametry~\cite{benedetti2019vqc,mitarai_quantum_2018}.
W rozdziałach~\ref{chapter:learned_encoding} i~\ref{chapter:classification} obwód nie ma własnych parametrów: sieć klasyczna $\Theta_{\mathbf{w}}$ wyznacza dla każdej próbki kąty $\boldsymbol{\theta}=\Theta_{\mathbf{w}}(\mathbf{x})$ obwodu o stałej strukturze, a wyjście modelu powstaje z wyników pomiaru przez klasyczne przetwarzanie końcowe (rys.~\ref{fig:tb:ml:hybrid}).
Uczone są więc wagi $\mathbf{w}$ sieci klasycznej i, jeśli występują, parametry przetwarzania końcowego (splot końcowy w rozdziale~\ref{chapter:learned_encoding}), a obwód jest ustaloną warstwą, nieliniową względem $\boldsymbol{\theta}$.
Gdy celem jest odtworzenie wejścia, model jest autoenkoderem hybrydowym: koderem jest sieć $\Theta_{\mathbf{w}}$ wraz z obwodem, reprezentacją pośrednią stan kwantowy, a dekoderem pomiar i przetwarzanie końcowe.

\begin{figure}[htb]
\centering
% Schemat hybrydowego modelu kwantowo-klasycznego z obwodem o stałej strukturze.
% Rysunek własny (PHD-TS-17); wstawiany przez \input w chapters/theoretical_background_06_ML.tex.
\begin{tikzpicture}[
  font=\small,
  >={Stealth[length=2mm]},
  box/.style={draw, rounded corners=2pt, align=center, minimum height=1.0cm, inner sep=3pt},
  flow/.style={->, thick},
  lbl/.style={font=\footnotesize, align=center},
]
% wejście i sieć klasyczna
\node[box, fill=gray!15, minimum width=0.8cm] (in) at (0,0) {$\mathbf{x}$};
\node[box, fill=blue!15] (net) at (1.9,0) {sieć\\klasyczna $\Theta_{\mathbf{w}}$};
\draw[flow] (in) -- (net);

% obwód o stałej strukturze
\foreach \y in {0.45,0,-0.45} {
  \node[font=\scriptsize] at (4.0,\y) {$\ket{0}$};
  \draw (4.25,\y) -- (7.35,\y);
}
\draw[fill=white] (4.6,-0.65) rectangle (5.6,0.65);
\node at (5.1,0) {$U(\boldsymbol{\theta})$};
\draw[dashed, fill=red!8] (5.85,-0.65) rectangle (6.45,0.65);
\node at (6.15,0) {$\channel$};
\foreach \y in {0.45,0,-0.45} {
  \draw[fill=white] (6.75,{\y-0.15}) rectangle (7.15,{\y+0.15});
  \draw (6.82,{\y-0.08}) arc (160:20:0.14);
  \draw[->, >={Stealth[length=1mm]}] (6.95,{\y-0.1}) -- (7.08,{\y+0.11});
}
\draw[dashed, rounded corners=3pt, gray!70!black] (3.7,-0.95) rectangle (7.5,0.95);
\node[lbl, text=gray!30!black] at (5.6,-1.2) {obwód o stałej strukturze};

% kąty z sieci klasycznej
\draw[flow] (net.north) |- (5.1,1.4) -- node[right, pos=0.4] {$\boldsymbol{\theta}=\Theta_{\mathbf{w}}(\mathbf{x})$} (5.1,0.65);

% odczyt, przetwarzanie końcowe, strata
\node[box, fill=blue!15, font=\footnotesize] (post) at (11.1,0) {przetwarzanie\\końcowe};
\draw[flow] (7.5,0) -- node[above, font=\scriptsize] {$\tilde{\prob}(\ket{b})$, $\langle O\rangle_{\boldsymbol{\theta}}$} node[below, font=\scriptsize] {$\nshots$ pomiarów} (post.west);
\node[box, fill=orange!20] (loss) at (13.45,0) {$\ell(\hat{\mathbf{y}},\mathbf{y})$};
\draw[flow] (post) -- node[above, font=\scriptsize] {$\hat{\mathbf{y}}$} (loss);

% gradient
\draw[->, dashed, thick, gray!60!black] (loss.south) -- ++(0,-1.0) -| (net.south);
\node[lbl, text=gray!30!black] at (6.9,-2.0) {gradient $\partial\symloss/\partial\mathbf{w}$: różniczkowanie automatyczne (symulator)\\lub model zastępczy urządzenia (dostrajanie na urządzeniu)};
\end{tikzpicture}

\caption[Hybrydowy model kwantowo-klasyczny z obwodem o stałej strukturze]{Hybrydowy model kwantowo-klasyczny z obwodem o stałej strukturze.
Sieć klasyczna $\Theta_{\mathbf{w}}$ wyznacza kąty $\boldsymbol{\theta}$ obwodu $U(\boldsymbol{\theta})$; blok $\channel$ oznacza szum (model szumu w symulatorze albo szum urządzenia).
Opracowanie własne.}
\label{fig:tb:ml:hybrid}
\end{figure}

Pochodne wyjścia obwodu względem $\boldsymbol{\theta}$, potrzebne do gradientu straty, wyznacza się w symulatorze automatycznym różniczkowaniem, a na urządzeniu można je oszacować regułą przesunięcia parametru, z dwóch wykonań obwodu z przesuniętym kątem na każdą składową gradientu~\cite{mitarai_quantum_2018,schuld_evaluating_2019}.
W rozdziałach~\ref{chapter:learned_encoding} i~\ref{chapter:classification} reguły tej się nie stosuje: w dostrajaniu na urządzeniu, czyli dalszym uczeniu nauczonego wcześniej modelu, wynik obwodu pochodzi z pomiaru, a gradient z różniczkowalnego modelu zastępczego urządzenia w symulatorze~\cite{wright_deep_2022}.
Uczenie z modelem szumu polega na zastąpieniu w symulatorze ewolucji unitarnej kanałami kwantowymi (podrozdział~\ref{subsec:tb:noise:models}); tak nauczony model jest dopasowany do założonego modelu szumu, a nie do urządzenia, którego szum może mieć inną strukturę i zmieniać się między kalibracjami.
Ponadto uczenie prowadzi się zwykle na dokładnych prawdopodobieństwach, a na urządzeniu dostępne są tylko estymaty obarczone szumem próbkowania; oba źródła błędu są w rozdziałach~\ref{chapter:learned_encoding} i~\ref{chapter:classification} analizowane osobno.

\subsection{Grupy Liego w bramkach kwantowych i sieciach neuronowych}\label{subsec:tb:ml:lie}\label{subsec:tb:lie}

Bramki jednokubitowe po pominięciu fazy globalnej należą do grupy Liego $\mathrm{SU}(2)$ macierzy unitarnych o wyznaczniku 1~\cite{hall2015lie}; elementy jej algebry Liego $\liealg{su}(2)$, czyli przestrzeni stycznej do grupy w elemencie neutralnym, zapisuje się jako $-\mathrm{i}\gen$ z hermitowskim, bezśladowym generatorem $\gen$, a odwzorowanie (mapa) wykładnicze przeprowadza je w grupę, co daje postać bramki
\begin{equation}\label{eq:tb01:generator}
  U(\theta)=\exp\Bigl(-\mathrm{i}\tfrac{\theta}{2}\gen\Bigr),
\end{equation}
w której obrotom~\eqref{eq:tb01:rotation} odpowiada $\gen=\sigma_a$.
Postać ta pozwala modelować błąd bramki zaburzeniem generatora, $\gen\mapsto\gen+\genpert$: zaburzenie hermitowskie zachowuje unitarność i odpowiada błędowi koherentnemu (odchyleniu osi lub kąta obrotu), a niehermitowskie daje operator nieunitarny~\cite{raubitzek_quantum_noise} (podrozdział~\ref{par:lr:noise:generators}).

Symetrię danych można wbudować w architekturę sieci~\cite{bronstein2021geometric}: odwzorowanie $f$ jest \glsdisp{equivariance}{ekwiwariantne} względem grupy $G$ działającej na wejściach przez przekształcenia $T_g$, a na cechach przez $T'_g$, gdy $f(T_g\mathbf{x})=T'_g f(\mathbf{x})$ dla wszystkich $g\in G$ i $\mathbf{x}$.
W preprincie Raubitzka i in.~\cite{raubitzek_lie_group_manuscript}, omawianym w podrozdziale~\ref{par:lr:cls:lie}, grupa Liego nie jest symetrią danych, lecz strukturą narzuconą na przekształcenie cech wewnątrz sieci: cechy osadza się w wektorach $\mathbf{h}_i$, na przykład dzieląc je na bloki, a warstwa działania grupy przekształca każdy z nich macierzą z grupy $G$, np. $\mathrm{SO}(3)$, $\mathrm{SU}(2)$ lub $\mathrm{SU}(3)$:
\begin{equation}\label{eq:tb:ml:lie}
\mathbf{h}_i\mapsto g_i\mathbf{h}_i,\qquad g_i=\exp\Bigl(\sum_{a=1}^{\dim\liealg{g}}c_{i,a}B_a\Bigr)\in G,
\end{equation}
gdzie $B_a$ są macierzami bazy algebry $\liealg{g}$ (w preprincie macierze bazy grup $\mathrm{SU}$ są hermitowskie, z jawnym czynnikiem $-\mathrm{i}/2$), a uczonymi parametrami są współrzędne $c_{i,a}$.
Parametryzacja przez algebrę gwarantuje, że $g_i$ należy do grupy przy dowolnych współrzędnych, a warstwę, klasyczny odpowiednik bramki obrotu~\eqref{eq:tb01:generator}, uczy się zwykłą propagacją wsteczną~\cite{raubitzek_lie_group_manuscript}.
